\section{A uniform bound on the variational expression}
\label{rcut:sec:scalar}

This appendix supplies the variational gap required by the planted
comparison. We use the notation of Section~\ref{rcut:sec:comparison}:
\(h\sim N(\beta^2,\beta^2)\), brackets denote expectation under this law,
and \(m=\tanh h\), \(v=1-m^2\), \(w=1-m\). The bound below has one
numerical margin for the entire stated temperature interval and for all
admissible functions.

\begin{proposition}[Uniform scalar gap]
\label{rcut:prop:scalar-gap}
For every \(\beta\in[0.35,0.5]\), every real measurable
\(b\in L^2\), and every measurable \(r\geq0\) in \(L^2\) satisfying
\(\langle b^2+r^2\rangle=1\), the expression in
Equation~\eqref{rcut:eq:variational} satisfies
\[
 \mathcal V_\beta(b,r)\leq 0.9969.
\]
\end{proposition}

Every terminating decimal in the proof denotes an exact rational number.
We will apply Young's inequality with a positive field-dependent
multiplier to separate the square-root term into quadratic terms in
\(b\) and \(r\). The multiplier has two jobs: the total diagonal
coefficient must stay below \(.9789\), and its cost on \(b^2\) must
absorb the remaining quadratic form, with an excess of at most
\(.018\langle b^2\rangle\). Normalization will then give
\(.9789+.018=.9969\).
The functional reductions are proved here. The scalar enclosures on
which they rely, including quadrature and arithmetic errors, are proved
in Appendix~\ref{rcut:sec:scalar-certificate}.

\subsection{Identities and the positive kernel remainder}
\label{rcut:sec:scalar-identities}

We begin with two identities that simplify the scalar coefficients.
Conditioned on \(|h|=t\), the sign of the field has weights proportional
to \(e^t\) and \(e^{-t}\), and hence mean \(\tanh t\). Multiplying by
an even function of the field gives, for every \(q\) for which the
expectations exist,
\begin{equation}
 \langle m q(m^2)\rangle=\langle m^2 q(m^2)\rangle.
 \label{rcut:eq:scalar-sign-identity}
\end{equation}
In particular, \(\langle m\rangle=\langle m^2\rangle\) and
\(\langle mv\rangle=\langle m^2v\rangle\).
The second identity follows by Gaussian integration by parts with
\(v'=-2mv\):
\begin{equation}
 \Cov(h,v)=-2\beta^2\langle mv\rangle,
 \qquad
 \langle(h-\beta^2+2\beta^2m)v\rangle=0.
 \label{rcut:eq:scalar-integration-parts}
\end{equation}

Choose the following positive multiplier for Young's inequality.
Its associated cost \(k\) will combine with \(D\) in the diagonal estimate:
\begin{equation}
 \begin{split}
 \Lambda{}(h)&=0.45+0.25m(h)^2-0.06m(h),\\
 k(y)&=\frac{\beta^2}{4}
       \left\langle\frac{(1-y^2+v)^2}{\Lambda{}}\right\rangle,
       \qquad -1\leq y\leq1,\\
 a&=\langle v\rangle,
 \qquad l=\langle \Lambda{}\rangle=0.64-0.19a.
 \end{split}
 \label{rcut:eq:scalar-multiplier}
\end{equation}
Equation~\eqref{rcut:eq:scalar-sign-identity} gives the expression for
\(l\) in Equation~\eqref{rcut:eq:scalar-multiplier}; since \(a\leq1\),
we have \(l\geq0.45\). Completing the square also gives
\(\Lambda{}=0.4464+0.25(m-0.12)^2\geq0.4464\).
These positive bounds justify the divisions below. To bound the
kernel term, define
\[
 g=\frac{v}{w+3/4},\qquad G=\langle g\rangle,
 \qquad
 R=\left\langle\frac{(1+m)v}{2}-\frac32 g^2\right\rangle.
\]

\begin{lemma}[Kernel remainder]
\label{rcut:lem:scalar-kernel}
The quantity \(R\) is nonnegative. For every real \(b\in L^2\),
\begin{equation}
 \langle b(h)K_0(h,h_*)b(h_*)\rangle_{h,h_*}
 \leq \frac32\langle bg\rangle^2+R\langle b^2\rangle.
 \label{rcut:eq:scalar-kernel-bound}
\end{equation}
\end{lemma}

\begin{proof}
We separate one rank-one term from the kernel and bound the remainder
by its diagonal. Put \(q=(w-3/4)/(w+3/4)\). A finite field has
\(w>0\) and \(|q|<1\), so expansion as a geometric series gives
\[
 K_0(h,h_*)=\frac32 g(h)g(h_*)\sum_{j=0}^{\infty}q(h)^j q(h_*)^j.
\]
The \(j=0\) term will give \(\frac32\langle bg\rangle^2\).
For the remaining terms set \(\phi_j=\sqrt{3/2}\,gq^j\), \(j\geq1\).
Their diagonal sum is
\begin{equation}
 \rho:=\sum_{j\geq1}\phi_j^2
 =\frac{(1+m)v}{2}\,q^2
 =\frac{(1+m)v}{2}-\frac32g^2\geq0.
 \label{rcut:eq:scalar-residual-diagonal}
\end{equation}
The bound \((1+m)v=(1-m)(1+m)^2\leq32/27\) implies
\(\rho\leq16/27\). Moreover \(\langle\rho\rangle=R\), proving that
\(R\geq0\). To justify the series for an arbitrary signed
\(b\in L^2\), apply Cauchy--Schwarz:
\[
 \sum_{j\geq1}\langle|b\phi_j|\rangle^2
 \leq\langle b^2\rangle\sum_{j\geq1}\langle\phi_j^2\rangle
 =R\langle b^2\rangle.
\]
This proves absolute summability of the signed double-integral series,
so its sum equals \(\sum_{j\geq1}\langle b\phi_j\rangle^2\).
The same estimate bounds that sum by \(R\langle b^2\rangle\).
Restoring the \(j=0\) term proves the kernel bound.
\end{proof}

\subsection{The scalar data}
\label{rcut:sec:scalar-data}

Write \(s_f=\langle(f-\langle f\rangle)^2\rangle^{1/2}\) for the
standard deviation of a real function \(f\).
Lemma~\ref{rcut:lem:scalar-data} provides the numerical inputs for both
parts of the argument. Its five endpoint enclosures control the diagonal
coefficient. To treat the remaining quadratic form, we will write
\(b=\langle b\rangle+b_0\), with \(\langle b_0\rangle=0\), and
bound the constant, centered, and mixed terms using the moment table.
Appendix~\ref{rcut:sec:scalar-certificate} proves all these enclosures.

\begin{lemma}[Certified scalar enclosures]
\label{rcut:lem:scalar-data}
At \(\beta=1/2\), define
\[
 d(y)=\frac14\left\langle\frac{v^2}{2-y-m}\right\rangle,
 \qquad -1\leq y\leq1,
\]
where at \(y=1\) the integrand is interpreted as \(v(1+m)\).
Then
\begin{equation}
 \begin{array}{c|ccccc}
 \text{quantity}& k(0)&k(1)&d(0)&d'(0)&d(1)\\ \hline
 \text{lower bound}&.4276&.0898&.0925&.0522&.2288\\
 \text{upper bound}&.4278&.0900&.0927&.0524&.2290
 \end{array}
 \label{rcut:eq:scalar-five-enclosures}
\end{equation}
All bounds in this display are strict.
For every \(\beta\in[7/20,1/2]\), define
\begin{equation}
 F=(v+a)h+2\beta^2mv+
       \left(0.50-\frac{\beta^2a}{l}\right)v
       +0.12m+3\beta^2Gg.
 \label{rcut:eq:scalar-F}
\end{equation}
At the five exact values of \(\beta\) displayed in
Table~\ref{rcut:tab:scalar-moments}, each of the eight entries after the first
column differs from its corresponding expectation or standard deviation
under \(N(\beta^2,\beta^2)\) by less than \(0.001\).
\begin{table}[htbp]
 \centering
 \caption{Certified centers for the moment quantities, with radius
 \(0.001\); the values of \(\beta\) are exact.}
 \label{rcut:tab:scalar-moments}
 \begin{tabular}{c@{\quad}rrrrrrrr}
 \toprule
 \(\beta\)&\(a\)&\(G\)&\(s_v\)&\(s_g\)&
 \(\langle mv\rangle\)&\(R\)&\(s_F\)&\(s_{mv}\)\\
 \midrule
 .350&.890&.554&.128&.084&.081&.015&.659&.238\\
 .400&.861&.546&.156&.093&.096&.018&.730&.248\\
 .445&.832&.537&.179&.101&.107&.022&.786&.253\\
 .475&.813&.530&.195&.107&.114&.025&.821&.255\\
 .500&.796&.524&.207&.113&.120&.027&.847&.255\\
 \bottomrule
 \end{tabular}
\end{table}
\end{lemma}

The function \(F\) is the coefficient of the mixed term after centering
\(b\), as verified in Equation~\eqref{rcut:eq:scalar-full-expansion}.
Since \(\langle b_0\rangle=0\), its contribution is bounded by
\(s_F\lvert\langle b\rangle\rvert\|b_0\|_2\).

\begin{lemma}[Diagonal envelope]
\label{rcut:lem:scalar-envelope}
For \(0.35\leq\beta\leq0.5\) and every real field \(h\),
\begin{equation}
 \Lambda{}(h)+D(h)+k(m(h))\leq0.9789.
 \label{rcut:eq:scalar-envelope}
\end{equation}
\end{lemma}

\begin{proof}
We first reduce the temperature interval to its upper endpoint and then
bound the dependence on the field. The Gaussian field has density
\[
 p_\beta(h)=\frac1{\sqrt{2\pi}\beta}
   \exp\left(-\frac{\beta^2}{2}+h-\frac{h^2}{2\beta^2}\right).
\]
The logarithmic derivative of \(\beta^2p_\beta(h)\) is
\(\beta^{-1}-\beta+h^2\beta^{-3}>0\) for \(\beta<1\).
Thus, for fixed \(y\), positivity of the integrands bounds the
corresponding integrals in \(D\) and \(k\) by those at \(\beta=1/2\).
At that endpoint, \(D(h)=d(m(h))\).

It remains to bound the endpoint expression uniformly in
\(-1\leq y\leq1\). Convexity of \(k\) as a function of \(y^2\)
gives \(k(y)\leq k(0)+(k(1)-k(0))y^2\). To treat \(d\), take
\(A=2-m>1\) and use the identity
\[
 \frac1{A-y}=\frac1A+\frac{y}{A^2}
             +\frac{y^2}{A^2(A-y)}
\]
which, together with \(A-y\geq A-1\), gives, for every \(-1\leq y\leq1\),
\[
 d(y)\leq d(0)+d'(0)y+
               y^2\bigl(d(1)-d(0)-d'(0)\bigr).
\]
After adding \(\Lambda{}\), the quadratic bounds give the following
upper bound for \(\Lambda{}+d+k\):
\begin{equation}
 \max\{.45+d(0)+k(0),\ .70+k(1)+d(1)-d'(0)\}
       +|d'(0)-.06|.
 \label{rcut:eq:scalar-envelope-arithmetic}
\end{equation}
We evaluate this bound using slightly enlarged intervals: replace the
enclosures in Equation~\eqref{rcut:eq:scalar-five-enclosures} by radius
\(.0003\) about the respective centers
\[
 .4277,\quad .0899,\quad .0926,\quad .0523,\quad .2289.
\]
The two entries of the maximum are then at most \(.9709\) and
\(.9674\), and the absolute-value term is at most \(.0080\).
Their sum proves Equation~\eqref{rcut:eq:scalar-envelope}.
\end{proof}

\subsection{Young's inequality and the centered coherent form}
\label{rcut:sec:scalar-reduction}

The diagonal envelope is now available. We next convert the radical
in the variational expression into quadratic terms, retaining its
negative square produced by expanding the variance. Set \(A=\langle b\rangle\) and
\(B=\langle bv\rangle\). Apply
\(\beta r\sqrt{T}\leq \Lambda{}r^2+\beta^2T/(4\Lambda{})\) to the last
term of Equation~\eqref{rcut:eq:variational}. Interchanging the two
expectations in the resulting quadratic terms gives
\[
 \langle(\Lambda{}+k(m))r^2\rangle
 \quad\hbox{and}\quad
 \langle k(m)b^2\rangle
   -\frac{\beta^2}{4}\left\langle\frac{(B+vA)^2}{\Lambda{}}\right\rangle.
\]
Weighted Cauchy--Schwarz bounds the negative term from above by
\(-\beta^2(B+aA)^2/(4l)\). Collect it with the products of expectations
and the kernel term, defining the coherent form
\begin{equation}
 \begin{split}
 \mathcal C_\beta(b)={}&\langle bh\rangle B
 +A\langle b(h-\beta^2+2\beta^2m)v\rangle\\
 &+\beta^2\langle bK_0b_*\rangle
 -\frac{\beta^2}{4l}(B+aA)^2,
 \end{split}
 \label{rcut:eq:scalar-coherent}
\end{equation}
and
\begin{equation}
 \mathcal V_\beta(b,r)
 \leq\langle(D+k(m))(b^2+r^2)\rangle
       +\langle \Lambda{}r^2\rangle+\mathcal C_\beta(b).
 \label{rcut:eq:scalar-young}
\end{equation}

\begin{lemma}[Coherent form bound]
\label{rcut:lem:scalar-coherent}
For every \(0.35\leq\beta\leq0.5\) and every real \(b\in L^2\),
\begin{equation}
 \mathcal C_\beta(b)\leq\langle(\Lambda{}+.018)b^2\rangle.
 \label{rcut:eq:scalar-coherent-bound}
\end{equation}
\end{lemma}

\begin{proof}
We reduce the inequality to a quadratic form in the constant and
centered parts of \(b\). Write \(b=A+b_0\), with
\(\langle b_0\rangle=0\), and set \(H_0=\langle b_0h\rangle\),
\(B_0=\langle b_0v\rangle\), \(G_0=\langle b_0g\rangle\).
Replace the kernel term by Lemma~\ref{rcut:lem:scalar-kernel}, and call
the resulting upper form \(\mathcal C_\beta^+\).
Equation~\eqref{rcut:eq:scalar-integration-parts} makes its purely constant
part \(\beta^2(a+\tfrac32G^2+R-a^2/l)A^2\). Define the two slacks
\begin{equation}
 \begin{split}
 U_0&=l+.018+\beta^2\left(\frac{a^2}{l}-a-\frac32G^2\right)
                -\beta^2R,\\
 U_1&=.4644-\frac{\beta s_v}{2}+\beta^2\langle mv\rangle
                -\beta^2\left(\frac32s_g^2+R\right).
 \end{split}
 \label{rcut:eq:scalar-slacks}
\end{equation}
Expanding the upper form around \(A+b_0\) and subtracting the desired
upper bound gives
\begin{equation}
 \begin{split}
 \mathcal C_\beta^+(b)-\langle(\Lambda{}+.018)b^2\rangle
 ={}&-U_0A^2+A\langle b_0F\rangle+H_0B_0
       +\frac32\beta^2G_0^2\\
 &+\beta^2R\langle b_0^2\rangle
       -\frac{\beta^2B_0^2}{4l}
       -\langle(\Lambda{}+.018)b_0^2\rangle.
 \end{split}
 \label{rcut:eq:scalar-full-expansion}
\end{equation}
The coefficient of the mixed term can be checked directly.
The first two terms in Equation~\eqref{rcut:eq:scalar-coherent} give
\((v+a)h+2\beta^2mv\); the negative square gives
\(-\beta^2av/l\); and the kernel bound gives \(3\beta^2Gg\).
Subtracting \(2\langle \Lambda{}b_0\rangle\) contributes
\(.50v+.12m\), up to a constant whose pairing with \(b_0\) vanishes.
These contributions sum to \(F\) from Equation~\eqref{rcut:eq:scalar-F}.

We next bound the terms quadratic in \(b_0\). On the centered space,
\[
 H_0B_0\leq\frac{\beta s_v+\Cov(h,v)}2\langle b_0^2\rangle
 =\left(\frac{\beta s_v}{2}-\beta^2\langle mv\rangle\right)
       \langle b_0^2\rangle.
\]
For this inequality, the symmetric rank-two operator associated with
\(u=h-\langle h\rangle\) and \(z=v-a\) has largest eigenvalue
\((\|u\|_2\|z\|_2+\langle uz\rangle)/2\).
Together with \(G_0^2\leq s_g^2\langle b_0^2\rangle\) and
\(\Lambda{}+.018\geq.4644\), this controls all centered quadratic
terms. Discarding the negative \(B_0^2\) term from
Equation~\eqref{rcut:eq:scalar-full-expansion} yields
\begin{equation}
 \mathcal C_\beta(b)-\langle(\Lambda{}+.018)b^2\rangle
 \leq-U_0A^2-U_1\langle b_0^2\rangle+A\langle b_0F\rangle.
 \label{rcut:eq:scalar-two-dimensional}
\end{equation}
By Cauchy--Schwarz, the mixed term is at most
\(s_F|A|\|b_0\|_2\). Thus positivity of \(U_0,U_1\) and the bound
\(s_F<2\sqrt{U_0U_1}\) will finish the proof. We verify these
inequalities on the entire temperature interval, using the full
radius \(.001\) of Table~\ref{rcut:tab:scalar-moments}.

First we establish the monotonicities needed between table entries.
The density of \(|h|\) is proportional to
\(\beta^{-1}e^{-\beta^2/2-h^2/(2\beta^2)}\cosh h\) on \(h\geq0\);
its ratio at a larger and a smaller \(\beta\) is increasing in \(h\).
To see the implication for expectations, let \(\ell\) be this ratio
and \(f\) a decreasing function. For independent \(X,X'\) under the
smaller-\(\beta\) law,
\[
 \Cov(f(X),\ell(X))
 =\frac12\E[(f(X)-f(X'))(\ell(X)-\ell(X'))]\leq0.
\]
Since \(\E\ell(X)=1\), the expectation of \(f\) decreases as
\(\beta\) increases. In particular \(a\) decreases. For \(G\),
average \(g\) over the conditional signs and write
\(u=\tanh^2|h|\); the resulting function is
\[
 \frac{(1-u)(1.75+u)}{1.75^2-u}.
\]
Its derivative,
\(4(64u^2-392u-35)/(16u-49)^2<0\) for \(0\leq u\leq1\), proves
that \(G\) decreases as well. The pointwise increase of
\(\beta^2p_\beta\) from Lemma~\ref{rcut:lem:scalar-envelope} gives the
other needed monotonicities: both
\(\beta^2\langle mv\rangle=\beta^2\langle m^2v\rangle\) and
\(\beta^2R\) increase. For \(\beta^2s_v^2\) and \(\beta^2s_g^2\),
write each weighted variance as the infimum over constant centers of
the integral of a nonnegative squared function against
\(\beta^2p_\beta\). Each such integral increases, so the infimum does too.

Now take successive temperatures \(j<z\) from
Table~\ref{rcut:tab:scalar-moments}. Enclose \(a\) on this interval
by taking \(a_-\) to be the center at \(z\) minus \(.001\), and
\(a_+\) the center at \(j\) plus \(.001\). Define \(G_-,G_+\)
in the same way, and set \(l_-=.64-.19a_+\),
\(l_+=.64-.19a_-\). For any endpoint quantity \(X_t\), write
\(X_t^+\) and \(X_t^-\) for its table center plus and minus
\(.001\). The monotonicities above give the following lower bounds
for the slacks throughout \([j,z]\):
\begin{equation}
 \begin{split}
 \underline U_0={}&l_-+.018
   +j^2\left(\frac{a_-^2}{l_+}-a_--\frac32G_+^2\right)-z^2R_z^+,\\
 \underline U_1={}&.4644-\frac{z(s_v)_z^+}{2}
   +j^2\langle mv\rangle_j^-
   -z^2\left(\frac32((s_g)_z^+)^2+R_z^+\right).
 \end{split}
 \label{rcut:eq:scalar-interval-slacks}
\end{equation}
For the first bound, use that \(a^2/l-a\) increases in \(a\) and
decreases in \(l\) on \(.795\leq a\leq.891\),
\(.45\leq l\leq.489\). The parenthesized coefficient in that line
is positive on each interval; its explicit lower bounds are given below.

To control the mixed coefficient, compare its standard deviation to
the value at the right endpoint. For \(j\leq\beta\leq z\), the
density ratio satisfies
\[
 \sup_h\frac{p_\beta(h)}{p_z(h)}
 =\frac z\beta\exp\left(\frac{z^2-\beta^2}{2}\right)
 \leq\frac zj\exp\left(\frac{z^2-j^2}{2}\right).
\]
Choose a constant center before changing measure. This bounds the
standard deviation of any function by its standard deviation under
\(p_z\) times the square root of the displayed factor. Next compare
the coefficients of \(F\) at \(\beta\) and \(z\), and apply the
triangle inequality for the centered \(L^2(p_z)\) seminorm. With
\(s_h=z\) under \(p_z\), this gives
\begin{equation}
 \begin{split}
 \overline s_F={}&
 \sqrt{\frac zj\exp\left(\frac{z^2-j^2}{2}\right)}
 \biggl[(s_F)_z^++(a_+-a_-)z
       +2(z^2-j^2)(s_{mv})_z^+\\
 &\hspace{15mm}
       +\left(\frac{z^2a_+}{l_-}-\frac{j^2a_-}{l_+}\right)(s_v)_z^+
       +3(z^2G_+-j^2G_-)(s_g)_z^+\biggr].
 \end{split}
 \label{rcut:eq:scalar-interval-cross}
\end{equation}
Substituting the full enclosures from
Table~\ref{rcut:tab:scalar-moments} gives the interval bounds in
Table~\ref{rcut:tab:scalar-intervals}. Appendix~\ref{rcut:sec:scalar-certificate}
specifies their directed-arithmetic verification. In the same interval
order, the parenthesized coefficients in the first line of
Equation~\eqref{rcut:eq:scalar-interval-slacks} exceed
\(.229787,.152558,.111290,.074675\), confirming the positivity used
above.

\begin{table}[htbp]
 \centering
 \caption{Strict lower bounds in the middle three columns and strict
 upper bounds in the last column, using the full radius \(.001\)
 in Table~\ref{rcut:tab:scalar-moments}.}
 \label{rcut:tab:scalar-intervals}
 \begin{tabular}{ccccc}
 \toprule
 \([j,z]\)&\(\underline U_0>\)&\(\underline U_1>\)&
 \(2\sqrt{\underline U_0\underline U_1}>\)&\(\overline s_F<\)\\
 \midrule
 \([.350,.400]\)&.513818&.437638&.948403&.842935\\
 \([.400,.445]\)&.514074&.431904&.942404&.896470\\
 \([.445,.475]\)&.515902&.429026&.940926&.897871\\
 \([.475,.500]\)&.513188&.426022&.935157&.914504\\
 \bottomrule
 \end{tabular}
\end{table}

The table therefore provides uniform margins
\(U_0>.513\), \(U_1>.426\), \(s_F<.920\), and
\(2\sqrt{U_0U_1}>.930\). These margins absorb the mixed term:
\[
 s_F|A|\|b_0\|_2
 \leq\frac{92}{93}\bigl(U_0A^2+U_1\|b_0\|_2^2\bigr).
\]
Substitution in Equation~\eqref{rcut:eq:scalar-two-dimensional} makes
its right side at most \(-(.426/93)\|b\|_2^2\), and in particular
nonpositive. This proves Equation~\eqref{rcut:eq:scalar-coherent-bound}
for every \(b\in L^2\).
\end{proof}

\begin{proof}[Proof of Proposition~\ref{rcut:prop:scalar-gap}]
We can now assemble the two bounds. Insert the coherent-form estimate
of Lemma~\ref{rcut:lem:scalar-coherent} into
Equation~\eqref{rcut:eq:scalar-young}, and then apply the diagonal
envelope of Lemma~\ref{rcut:lem:scalar-envelope}:
\[
 \mathcal V_\beta(b,r)
 \leq\langle(D+k(m)+\Lambda{})(b^2+r^2)\rangle+.018\langle b^2\rangle
 \leq.9789+.018=.9969.
\]
The argument applies to the full function class in the proposition.
The coefficients \(h,hv\) belong to \(L^2\), all other coefficient
functions are bounded, and Lemma~\ref{rcut:lem:scalar-kernel} already
justifies the kernel integral for signed, unbounded \(b\).
The quantity under the last radical in
Equation~\eqref{rcut:eq:variational} is at most
\(4\langle b^2+r^2\rangle=4\), making its pairing with \(r\) finite.
Thus Fubini for the nonnegative quadratic integrals and every use of
Cauchy--Schwarz are valid without an infinity-norm bound on either
function. The estimate consequently also covers concentration on
arbitrarily rare field sets.
\end{proof}

